%% kapa: exam=1 round=26 year=2015 period=1 form=A subject=civil_law subject_ko=민법 seq=38 number=38
%% source: Q-Net 감정평가사 기출문제 (cst00309) artlSeq=5138001 file=extracted/1008/5138001/제26회_감정평가사 1차_기출/제26회_감정평가사 1차_기출/1교시_감평_A형.pdf
%% conversion: kordoc 4.24.51 -> tools/gen1.py (auto); review=proofread-sonnet; flags=none
\begin{question}{38}
전세권에 관한 설명으로 옳지 않은 것은? (다툼이 있으면 판례에 따름)
\begin{choices}
  \item 전세권자는 전세권설정계약에 다른 약정이 없는 한 전세권설정자의 동의 없이 전전세를 할 수 있다.
  \item 전세권이 성립한 후 전세목적물의 소유권이 이전된 경우, 전세권은 전세권자와 목적물의 소유권을 취득한 신 소유자 사이에서 동일한 내용으로 존속한다.
  \item 전세권이 갱신없이 존속기간이 만료되면 그 용익 물권적 권능은 전세권설정등기의 말소 없이도 소멸한다.
  \item 전세권자는 전세권설정자에게 목적물의 현상을 유지하기 위하여 지출한 필요비의 상환을 청구할 수 있다.
  \item 건물의 일부에 대하여 전세권이 설정되어 있는 경우, 그 전세권의 목적이 된 부분이 구조상 또는 이용상 독립성이 없어 그 부분만의 경매신청이 불가능하다고 하더라도, 이를 이유로 전세권의 목적물이 아닌 나머지 건물부분에 대하여 그 전세권에 기한 경매신청을 할 수 없다.
\end{choices}
\end{question}
\answer{4}
